\section{Metastability and Kinetic Trapping}

Equilibrium diagrams describe the thermodynamic destination, whereas experiments traverse a kinetic landscape. In oxide glasses, structural relaxation changes the population of nuclei and can produce distinct surface- and volume-crystallization pathways \cite{mckenzie2020nucleation,mckenzie2021nucleation,schmelzer2020effects,waurischk2020crystal}. Barium-silicate network studies connect hetero- and homo-connectivity to the ease of nucleation, providing a mechanistic reason why two glasses with similar bulk composition can crystallize differently \cite{moulton2021a,moulton2022a}.

Low-temperature intermediates and mechanochemical activation can template a metastable polymorph before the equilibrium phase becomes kinetically accessible. Layered-oxide studies identify such intermediates during solid-state synthesis, while mechanochemical work shows that metastable phases can be retained by milling history and short anneals \cite{bai2020kinetic,rana2021insights}. Recent non-equilibrium synthesis studies generalize this principle to oxide discovery: deliberate control of precursor and thermal path can select products outside the equilibrium sequence \cite{pitcher2024non,zeng2024selective}.

Interfacial melts create a second trapping mechanism. A small liquid fraction can accelerate transport, alter local stoichiometry, or stabilize a phase that would not nucleate in a dry solid mixture. Thermodynamic analyses of interfacial-melt stability and impurity-minimizing reaction networks make this coupling explicit \cite{zhang2026interfacial,jang2025identification,dhar2026thermodynamic}. In a Ba--Y--Zn--Si--O experiment, a Pt or Au crucible, cover powder, or flux may therefore be part of the reaction network rather than a passive container \cite{thieme2022noble,abdeyazdan2024phase}.

Quenching preserves a snapshot of the high-temperature assemblage but can also freeze a glass or metastable polymorph. Time--temperature studies show that finite holds and cooling rates change the fraction of kinetic by-products, so a single furnace schedule cannot define an equilibrium phase \cite{wang2024optimal,wang2024optimal,walters2026constructing}. The discriminating experiment is a matrix of identical compositions equilibrated for increasing hold times, sampled by rapid quench and slow cool, and analyzed by XRD plus EPMA and thermal methods. Agreement across time and cooling histories supports an equilibrium assignment; systematic dependence indicates kinetic control.
